% [原文 p.163]
\chapter[K 理论与实结构]{K 理论与实结构\protect\footnote{M.~F.~Atiyah, `K-theory and Reality', \textit{Q. J. Math.} \textbf{17}(1)(1966) 367--386.}}

\para{引言}

复向量丛的 K 理论[2],[5]有如下许多变体与精细化：

\begin{enumerate}[label=(\arabic*),leftmargin=2.2em,itemsep=2pt]
\item 实数上的向量丛的 K 理论 $KO$；
\item 自共轭丛的 K 理论 $KC$[1] 或 $KSC$[7]；
\item $G$ 空间上的 $G$ 向量丛的 K 理论 $K_G$[6]。
\end{enumerate}

本论文引入一种在某种意义上把这三者混合起来的新 K 理论 $KR$。定义 $KR$ 的动机之一是与实代数几何的类似，另一部分则来自实椭圆算子的理论。实际上，要彻底处理实椭圆算子的指数，这个 $KR$ 理论是本质的。另一方面，从纯粹拓扑的观点看，$KR$ 理论也有若干优点：把 $KR$ 作为中心理论、把其他 K 群视为 $KR$ 的派生，在许多场合是重要的。使用 $KR$ 理论，如[3]所证明的，可以把 K 理论的周期性作为本质的出发点，给出 $KO$ 理论周期性的漂亮证明。这个证明是本论文的主要目的之一。在此证明的过程中，自共轭理论 $KSC$ 与其他各种理论之间的种种完全序列会以自然的形式出现。这与 Anderson 的硕士论文[1]有相当大的重叠，但按本论文这种更通透的新观点，各种理论之间的关系会更容易看清。

最近，Karoubi 展开了关于带有对合的范畴的抽象 K 理论[8]。本论文的理论包含在这个抽象理论之中，但[8]并没有展开 $KR$ 特有的性质，
% [原文 p.164]
也不是用来简化 $KO$ 理论的周期性的。C.1 节给出 $KR$ 的定义与基本性质。C.2 节论述 $KR$ 的周期性定理以及它作为一般上同调理论的性质。C.3 节处理各种派生的理论：引入以某个固定空间为系数的 $KR$ 理论，并以关于 $KO$ 的周期性定理收尾。C.4 节按与[4]同样的方式简略论述 $KR$ 与 Clifford 代数的关系，特别证明 C.3 节所用的引理。C.5 节简略论述实椭圆算子的拓扑研究中 $KR$ 理论的重要性。

本论文本质上是与 I.~M.~Singer 关于指数定理的合作研究的副产品。由于这一结果具有独立于指数定理的拓扑意义，我们认为单独发表更好。

\section{实空间的范畴}
\takeaway{对合空间（实空间）以及沿对合的反线性共轭，是本节的全部语言；实向量丛、实空间之间的映射以及同伦都在这一范畴中定义，$KR$ 的载体由此就位。}

对合空间是指拓扑空间 $X$ 与周期 2 的同胚 $\tau\colon X\to X$（即 $\tau^2$ 为恒等映射）组成的对 $(X,\tau)$。对合 $\tau$ 视为 $X$ 的结构的一部分，在没有混淆之虞时常常省略。把 $\Z/2$ 视为 2 阶群时，对合空间就是[6]意义下的 $\Z/2$ 空间。另一个更具启发性的术语是：把对合空间称为实空间。这来自代数几何。实际上，若 $X$ 是实代数流形的复点集合，则对合由复共轭给出，$X$ 具有本论文意义下实空间的自然结构。注意这个对合的不动点恰好是流形 $X$ 的实点。仿照这个例子，把对合 $\tau$ 像复共轭那样写作
\[
\tau(X)=\bar x.
\]
% [Tutor Note: 原书 p.164 此式左端印作大写的 $\tau(X)$，右端却是 $\bar x$；按下文“把对合像复共轭那样写”的用法，应为 $\tau(x)=\bar x$，疑为排版笔误。]

“实空间”$X$ 上的“实向量丛”是指 $X$ 上的复向量丛 $E$，它同时是实空间，且满足下列条件：

\begin{enumerate}[label=(\roman*),leftmargin=2.4em,itemsep=2pt]
\item 投影 $E\to X$ 是实的（即与 $E$ 上以及 $X$ 上的对合可交换）；
\item 映射 $E_x\to E_{\bar x}$ 是反线性的，即下图可交换。
\end{enumerate}

% [原文 p.165]
\[
\begin{array}{ccc}
\C\times E_x & \longrightarrow & E_x\\
\downarrow & & \downarrow\\
\C\times E_{\bar x} & \longrightarrow & E_{\bar x}
\end{array}
\]
其中竖向的箭头表示对合，$\C$ 上赋有标准的实结构（$\tau(z)=\bar z$）。

重要的是，要注意（如此定义的）实空间范畴中的向量丛与 $\Z/2$ 空间范畴中的复向量丛之间的区别。在后一种情形中，映射
\[
E_x\to E_{\tau(x)}
\]
是复线性的。又若 $E$ 是 $\Z/2$ 空间范畴中的实向量丛，则其复化会因把
\[
E_x\to E_{\tau(x)}
\]
线性地扩张还是反线性地扩张，而得到两种不同的结构：线性扩张得到 $\Z/2$ 空间范畴中的复向量丛，反线性扩张则得到实空间范畴中的向量丛。

% [Tutor Note: 这段话是理解全篇的关键：同一个复向量丛 $E$ 配上“线性”的对合就落在 $\Z/2$ 空间范畴里，配上“反线性”的对合才落在实空间范畴里；$KR$ 与 $K_{\Z/2}$ 的差别正在于此。]

在 $X$ 上对合的不动点（也称为 $X$ 的实点）处，$E$ 上的对合给出反线性映射
\[
\tau_x\colon E_x\to E_x,
\]
且满足 $\tau_x^2=1$。这意味着 $E_x$ 是由一个实向量空间以自然方式复化而来，即 $\tau_x$ 的 $+1$ 特征子空间（$E_x$ 的实点）的复化。特别地，若 $X$ 上的对合是恒等的、$X$ 的一切点都是实点，则 $X$ 上通常意义下的实向量丛的范畴 $\cE(X)$ 与按本论文定义的意义下 $X$ 上实向量丛的范畴 $\cF(X)$\footnote{$\cF(X)$ 中的态射留到后面再定义。} 之间存在自然的等价。为此，只需（给 $\C$ 以标准的实结构而）用 $E\mapsto E\otimes_\R\C$ 定义 $\cE(X)\to\cF(X)$，并（以 $F_\R$ 表示 $F$ 的实点集合）用 $F\mapsto F_\R$ 定义 $\cF(X)\to\cE(X)$ 即可。

% [原文 p.166]
经由这个范畴等价，在实空间上使用“实向量丛”这一术语就得到了正当化。也就是说，“实空间上的实向量丛”是通常空间上实向量丛的自然推广。

若 $E$ 是实空间 $X$ 上的实向量丛，则截面空间 $\Gamma(E)$ 是带有反线性对合的复向量空间。也就是说，若 $s\in\Gamma(E)$，则 $\bar s$ 由
\[
\bar s(x)=\overline{s(\bar x)}
\]
定义。因此 $\Gamma(E)$ 具有实结构，即 $\Gamma(E)$ 是实向量空间 $\Gamma(E)_\R$ 的复化。

若 $E, F$ 是实空间 $X$ 上的实向量丛，则射 $\phi\colon E\to F$ 定义为与对合交换的复向量丛同态，即满足
\[
\phi(\bar e)=\overline{\phi(e)}\qquad(e\in E)
\]
的复向量丛同态。$E\otimes_\C F$ 与 $\Hom_\C(E,F)$ 具有自然的实向量丛结构。例如，若 $\phi_x\in\Hom_\C(E_{\bar x},F_{\bar x})$，则 $\overline{\phi_x}\in\Hom_\C(E_{\bar x},F_{\bar x})$ 由
\[
\overline{\phi_x}(u)=\overline{\phi_x(\bar u)}\qquad(u\in E_{\bar x})
\]
定义。于是，射 $\phi\colon E\to F$ 恰好是 $\Hom_\C(E,F)$ 的实截面，即 $(\Gamma\Hom_\C(E,F))_\R$ 的元素。

% [Tutor Note: 上面的 $\phi_x\in\Hom_\C(E_{\bar x},F_{\bar x})$ 与紧接的 $\overline{\phi_x}(u)=\overline{\phi_x(\bar u)}$（$u\in E_{\bar x}$）不相容：右式要求 $\phi_x$ 定义在 $E_x$ 上，故第一个 $\Hom_\C$ 的纤维应为 $E_x,F_x$，页图上 $x$ 的上横线疑为误排。]

又若 $X$ 紧，则正如[3, §1]那样，可以知道实向量丛的同伦论性质。唯一值得注意的是这一点：$X$ 的实子空间 $Y$ 上的实截面总能扩张为 $X$ 上的实截面。事实上，若 $t$ 是扩张 $s$ 的任一截面，则 $\frac12(t+\bar t)$ 就是作为实截面的扩张。

再设 $X$ 为实代数流形（即实代数流形的复点）。如已所述，$X$ 以自然的方式 $X_{\mathrm{alg}}\mapsto X_{\mathrm{top}}$ 定义实拓扑空间。为这里的目的，不妨把实代数向量丛取为复代数向量丛 $\pi\colon E\to X$，使 $X, E, \pi$ 以及数乘 $\C\times E\to E$ 全都在 $\R$ 上有定义（即由实系数公式给出）。着眼于作为底层结构的拓扑结构，$E_{\mathrm{top}}$ 显然是实空间 $X_{\mathrm{top}}$ 上的
% [原文 p.167]
实向量丛。

作为具体例子，考虑 $(n-1)$ 维复射影空间 $X=P(\C^n)$。$P(\C^n)$ 上的标准线丛 $H$ 是实代数向量丛。事实上，$H$ 由向量丛的完全序列
\[
0\to E\to X\times\C^n\to H\to0
\]
定义，其中 $E\subset X\times\C^n$ 由满足条件
\[
\sum u_iz_i=0
\]
的一切对 $((z),u)\in X\times\C^n$ 组成；这个等式的系数是实数，所以 $E$ 是实向量丛，由此又可推出 $H$ 也是实向量丛。因此，$H$ 定义了实空间 $P(\C^n)$ 上的实向量丛。

再举一例，考虑下列仿射二次曲面：
\[
\sum_{i=1}^n z_i^2+1=0.
\]
它是仿射的，故实向量丛可由仿射环 $A_+=\R[z_1,\dots,z_n]/(\sum z_i^2+1)$ 上的投射模来定义。这个二次曲面与虚轴的相交部分是球面
\[
\sum_{i=1}^n y_i^2=1,
\]
而对合正是对跖映射 $y\mapsto-y$ 本身。因此，仿射环 $A_+$ 上的投射模定义 $S^{n-1}$ 上一个以该对跖映射为对合的实向量丛。把上述二次曲面换成另一个二次曲面
\[
\sum z_i^2-1=0,
\]
它与实数平面的相交部分是球面，且具有平凡的对合，于是
\[
A_-=\frac{\R[z_1,\dots,z_n]}{(\sum z_i^2-1)}
\]
上的投射模定义 $S^{n-1}$ 上一个具有平凡对合（因而就是通常意义下的实向量丛）的实向量丛。这个例子中 $S^{n-1}$ 的意义在于，实空
% [原文 p.168]
间范畴中它是该二次曲面的形变收缩核（即保持对合的收缩）。

用 $KR(X)$ 表示实空间 $X$ 上实向量丛的 Grothendieck 群。限制到 $X$ 的实点上，就得到同构映射
\[
KR(X)\to KR(X_\R)\cong KO(X_\R).
\]
特别地，若 $X=X_\R$，则
\[
KR(X)\cong KO(X).
\]
例如取 $X=P(\C^n)$，则 $X_\R=P(\R^n)$，因此得到限制同构映射
\[
KR(P(\C^n))\to KR(P(\R^n))=KO(P(\R^n)).
\]
注意，在这个同态映射下 $[H]$ 的像正是 $P(\R^n)$ 上的标准实 Hopf 丛。

一如既往，$KR(X)$ 借助张量积成为环。

忘记 $X$ 上的对合，就得到自然同态映射
\[
c\colon KR(X)\to K(X).
\]
若 $X=X_\R$，这恰好是复化。另一方面，若 $E$ 是 $X$ 上的复向量丛，则 $E\oplus\tau^*\bar E$ 具有自然的实结构，因此得到同态映射
\[
\tau\colon K(X)\to KR(X).
\]
若 $X=X_\R$，这恰好是“实化”，即考虑作为底层结构的实空间。

\section{周期性定理}
\takeaway{把 Bott 周期性证明中的每一步都搬到实空间范畴，就得到 $KR$ 的 8 周期性；本节还给出它的上同调论性质，并由此把 $KO$ 的周期性化为 $KR$ 的周期性。}

下面进入周期性定理。这里将仔细追随[3, §2]的证明
% [原文 p.169]
展开，并指出本论文的理论所需要的修改点。

若 $E$ 是实空间 $X$ 上的实向量丛，则 $E$ 的射影丛 $P(E)$ 也是实空间。此外，$P(E)$ 上的标准线丛 $H$ 是实线丛。此时，关于 $KR$ 的周期性定理如下。

\begin{theorem}{C.2.1}
  设 $L$ 为紧的实空间 $X$ 上的实线丛，$H$ 为实空间 $P(L\oplus1)$ 上的标准实线丛。此时，作为 $KR(X)$-代数，$KR(P(L\oplus1))$ 由 $H$ 生成，并服从唯一的关系式
\[
([H]-1)([L][H]-1)=0.
\]
\end{theorem}

首先，选取 $L$ 的度量，使其在对合下不变。此时单位圆丛 $S$ 是实空间。由包含映射 $S\to L$ 定义的 $\pi^*(L)$ 的截面 $z$ 是“实”截面。因此，它的幂 $z^k$ 也是实截面。同构映射
\[
H^k\cong(1,z^{-k},L^{-k})\qquad([3,\ (2.5)])
\]
是实向量丛的同构映射。最后，我们说明：若 $f$ 是 $\Hom(\pi^*E^0,\pi^*E^\infty)$ 的实截面，则其傅里叶系数 $a_k$ 是 $\Hom(L^k\otimes E^0,E^\infty)$ 的实截面。事实上，
\begin{align*}
\overline{a_k(x)}&=\overline{a_k(\bar x)}=-\frac{1}{2\pi i}\int_{S_{\bar x}}\overline{f_{\bar x}z_{\bar x}^{-k-1}dz_{\bar x}}\\
&=\frac{1}{2\pi i}\int_{S_x}\overline{f_{\bar x}}\,(\overline{z_{\bar x}})^{-k-1}\,d\overline{z_{\bar x}}\quad(\text{因对合反转 }S\text{ 的定向})\\
&=\frac{1}{2\pi i}\int_{S_x}f_xz_x^{-k-1}dz_x\quad(\text{因 }f\text{ 与 }z\text{ 是实截面})\\
&=a_k(x)
\end{align*}
成立。也许想一想 $X$ 的实点处发生了什么会更容易理解。当 $x$ 是实点时，$f_x$ 为实截面这一条件成为
\[
f_x(\ee^{-\ii\theta})=\overline{f_x(\ee^{\ii\theta})},
\]
由此立刻推出傅里叶系数是实数。

% [原文 p.170]
[3, §3] 的线性化步骤只用到 $a_k$ 与 $z^k$，因此可知那里得到的同构映射全都是实同构映射。

若 $p$ 是实算子，则[3, §4]中的投影算子 $Q^0$ 与 $Q^\infty$ 也是实算子。事实上，
\begin{align*}
\overline{Q_x^0}&=\overline{Q_{\bar x}^0}=-\frac{1}{2\pi i}\int_{S_{\bar x}}\overline{p_{\bar x}^{-1}dp_{\bar x}}\\
&=\frac{1}{2\pi i}\int_{S_x}(\overline{p_{\bar x}})^{-1}d\overline{p_{\bar x}}\\
&=\frac{1}{2\pi i}\int_{S_x}p_x^{-1}dp_x\quad(\text{因 }p\text{ 是实算子}).
\end{align*}
$Q^\infty$ 的情形也相同。因此，向量丛 $V_n(E^0,p,E^\infty)$ 是实向量丛，而[3, (4.6)]是 $KR(P)$ 中的等式。剩下的只需照搬[3, §5]的证明即可。

至此，用空间对与悬垂来展开通常的一般上同调理论所需的准备已经就绪。不过，还会新增一个重要的特征。除基于具有平凡对合的 $\R$ 的通常悬垂之外，还可以考虑具有对合 $x\mapsto-x$ 的 $\R$。我们总给复数 $\C$ 赋以由复共轭给出的标准实结构，并把前者视为实轴 $\R\subset\C$、后者视为虚轴 $\ii\R\subset\C$，这样往往更方便。以下采用这样的记法：
\begin{gather*}
\R^{p,q}=\R^q\oplus\ii\R^p,\\
B^{p,q}=\R^{p,q}\text{ 的单位球体},\\
S^{p,q}=\R^{p,q}\text{ 的单位球面}.
\end{gather*}
注意 $\R^{p,p}\cong\C^p$。又注意按此记法 $S^{p,q}$ 的维数是 $p+q-1$。

空间对的 $KR$ 群 $KR(X,Y)$ 照常定义为 $\widetilde{KR}(X/Y)$，其中 $\widetilde{KR}$ 是限制到基点的核。此时把 $(p,q)$ 悬垂定义为
\[
KR^{p,q}(X,Y)=KR\bigl(X\times B^{p,q},\ X\times S^{p,q}\cup Y\times B^{p,q}\bigr).
\]
于是，通常的悬垂群 $KR^{-q}$ 由
% [原文 p.171]
\[
KR^{-q}=KR^{0,q}
\]
给出。于是，与[2]一样，得到实空间对 $(X,Y)$ 的完全序列
\begin{equation}
\cdots\to KR^{-1}(X)\to KR^{-1}(Y)\to KR(X,Y)\to KR(X)\to KR(Y)\tag{C.2.2}
\end{equation}
同样也可得到实空间三元组 $(X,Y,Z)$ 的完全序列。于是，对每个整数 $p\ge0$，得到三元组
\[
(X\times B^{p,q},\ X\times S^{p,0}\cup Y\times B^{p,0},\ X\times S^{p,0})
\]
的完全序列
\[
\cdots\to KR^{p,1}(X)\to KR^{p,1}(Y)\to KR^{p,0}(X,Y)\to KR^{p,0}(X)\to KR^{p,0}(Y).
\]

% [Tutor Note: 上面三元组的第一个空间页图作 $X\times B^{p,q}$，而同组的另两个空间与整个完全序列都只用到 $KR^{p,0}$、$KR^{p,1}$，此处应为 $X\times B^{p,0}$，疑为笔误。]

$KR(X)$ 的环结构以自然的方式扩张，使得有外积
\[
KR^{p,q}(X,Y)\otimes KR^{p',q'}(X',Y')\to KR^{p+p',q'+q'}(X'',Y''),
\]
其中 $X''=X\times X'$，$Y''=X\times Y'\cup X'\times Y$。把这些外积限制到对角部分，就定义了内积。

定理 C.2.1 可以照常重新表述。令
\[
b=[H]-1\in KR^{1,1}(\ast)=KR(B^{1,1},S^{1,1})=KR(P(\C^2)),
\]
并用 $\beta$ 表示由 $x\mapsto b\cdot x$ 给出的同态映射
\[
KR^{p,q}(X,Y)\to KR^{p+1,q+1}(X,Y).
\]
此时下列定理成立。

\begin{theorem}{C.2.3}
  $\beta\colon KR^{p,q}(X,Y)\to KR^{p+1,q+1}(X,Y)$ 是同构映射。
\end{theorem}

注意实空间对的完全序列与周期性定理的同构映射是相容的。因此，对 $p\ge0$ 定义
\[
KR^p(X,Y)=KR^{p,0}(X,Y).
\]
% [原文 p.172]
则可知关于 $(X,Y)$ 的完全序列 (C.2.2) 可以向两侧无限延伸。此外，还得到自然同构映射 $KR^{p,q}\cong KR^{p-q}$。

下面考虑[2, §2.7]中对 K 理论证明的一般 Thom 同构定理。回忆一下，其证明的核心步骤是这样进行的。

\begin{enumerate}[label=(\roman*),leftmargin=2.4em,itemsep=1pt]
\item 对线丛使用定理 C.2.1；
\item 对可分解的向量丛，用基于定理 C.2.1 的归纳法证明；
\item 对一般的向量丛，使用分裂原理。
\end{enumerate}

考察[2, §2.7]的证明可知，需要作本质修改的只有一点：向量丛局部平凡，因而局部可分解。实向量丛是按“具有实结构的向量丛”定义的，所以它作为通常空间范畴中的向量丛已被假定为局部平凡。需要证明的是，实向量丛在实空间范畴中也是局部平凡的。为此必须分两种情形讨论。

\begin{enumerate}[label=(\roman*),leftmargin=2.4em,itemsep=2pt]
\item $x\in X$ 是实点的情形。此时作为实空间有 $E_x\cong\C^n$。于是由 Tietze 扩张定理，存在 $x$ 的实邻域 $U$，使作为实空间有 $E|_U\cong U\times\C^n$。
\item $x\neq\bar x$ 的情形。取复同构映射 $E_x\cong\C^n$。它诱导出同构映射 $E_{\bar x}\cong\C^n$。于是取 $Y=\{x,\bar x\}$，得到实同构映射
\[
E|_Y\cong Y\times\C^n.
\]
由 Tietze 扩张定理，存在 $Y$ 的实邻域 $U$，使 $E|_U\cong U\times\C^n$。
\end{enumerate}

如此便得到下列定理。

\begin{theorem}{C.2.4}（Thom 同构定理）
  设 $E$ 为紧的实空间 $X$ 上的实向量丛。此时
% [原文 p.173]
\[
\phi\colon KR(X)\to\widetilde{KR}(X^E)
\]
是同构映射。其中 $\phi(x)=\lambda_E\cdot x$，而 $\lambda_E$ 是由 $E$ 的外代数确定的 $\widetilde{KR}(X^E)$ 的元素。
\end{theorem}

在[2, §2.7]的若干结果中，这里提及下列同构：
\begin{align*}
KR(X\times P(\C^n))&\cong KR(X)[t]/(t^n-1)\\
&\cong KR(X)\otimes_\Z K(P(\C^n)).
\end{align*}

Grassmann 流形与旗流形上 $KR$ 的计算，作为练习题留给读者。确定二次曲面上的 $KR$ 是更有趣的问题。这是因为其答案依赖于二次曲面的符号差。

下面用一项考察结束本节。考虑包含映射
\[
\R^{0,1}=\R\xrightarrow{i}\C=\R^{1,1}.
\]
这个包含映射诱导出同态映射
\[
\begin{array}{ccc}
KR^{1,1}(\ast) & \xrightarrow{i^*} & KR^{0,1}(\ast)\\
\Vert & & \Vert\\
\widetilde{KR}(P(\C^2)) & \longrightarrow & \widetilde{KR}(P(\R^2))
\end{array}
\]
由于 $i^*[H]$ 是 $P(\R^2)$ 上的实 Hopf 丛，$\eta=i^*(b)=i^*([H]-1)$ 是 $P(\R^2)$ 上的约化 Hopf 丛。

\section{系数理论}
\takeaway{固定一个实空间 $Y$，$X\mapsto KR(X\times Y)$ 给出以 $Y$ 为系数的上同调论；取 $Y=S^{p,0}$ 时，这一族理论与 $KO$ 的 8 周期性相互印证。}

设 $Y$ 为实空间。函子 $X\mapsto KR(X\times Y)$ 给出实空间上的新上同调理论。可以把这个理论称为以 $Y$ 为系数的 $KR$ 理论。取 $Y$ 为球面 $S^{p,0}$ 来看（其对合是对跖映射）。若理论 $F$ 有自然同构 $F\cong F^{-q}$，就说 $F$ 具有周期 $q$。于是得到下列命题。

% [原文 p.174]
\begin{proposition*}{C.3.1}
  以 $S^{p,0}$ 为系数的 $KR$ 理论的周期为
  \begin{gather*}
  p=1\ \text{时为}\ 2,\\
  p=2\ \text{时为}\ 4,\\
  p=4\ \text{时为}\ 8.
  \end{gather*}
\end{proposition*}

\begin{proof}
  当 $p=1,2,4$ 时，分别把 $\R^p$ 视为域 $\R$、$\C$、$\HH$。此时，对任意实空间 $X$，由 $\mu_p(x,s,u)=(x,s,su)$ 给出的映射
\[
\mu_p\colon X\times S^{p,0}\times\R^{0,p}\to X\times S^{p,0}\times\R^{p,0}
\]
是实同构映射，其中 $su$ 是 $\R$、$\C$、$\HH$ 中的乘法。于是这个映射诱导出同构映射
\[
\mu_p^*\colon KR^{p,0}(X\times S^{p,0})\to KR^{0,p}(X\times S^{p,0}).
\]
把 $X$ 换成它的悬垂，就得到同构映射
\[
\mu_p^*\colon KR^{p,q}(X\times S^{p,0})\to KR^{0,p+q}(X\times S^{p,0}).
\]
令 $q=p$，再用定理 C.2.1 给出的同构映射
\[
\beta^p\colon KR\to KR^{p,p},
\]
最终得到同构映射
\[
\begin{array}{ccc}
\mu_p^*\beta^p\colon KR(X\times S^{p,0}) & \longrightarrow & KR^{0,2p}(X\times S^{p,0})\\
& & \Vert\\
& & KR^{-2p}(X\times S^{p,0})
\end{array}
\]
\end{proof}

\begin{remark}{}
  $\mu^*$ 显然是 $KR(X)$-模同态。对 $\beta$ 也可说同样的话，因此由周期性定理的同构映射
\[
\gamma_p=\mu_p^*\beta^p\colon KR(X\times S^{p,0})\to KR^{-2p}(X\times S^{p,0})
\]
给出的映射，就是在同构映射 $KR(S^{p,0})\to KR^{-2p}(S^{p,0})$ 下乘以 $1$ 的像 $c_p$。这个元素 $c_p$ 由
\[
c_p=\gamma_p(1)=\mu^*(b^p\cdot1),\qquad 1\in KR(S^{p,0})
\]
给出。
\end{remark}

% 块尾跨页句：周期性定理の同型写像（p.174 末行起的 $\gamma_p$ 定义）——p.175 首句“は，同型写像 … における 1 の像 c_p による掛け算である．”已在本 chunk 写完；appC2 请从 p.175“任意の Y に対して，射影 X×Y→X は…”接续，不要重写该句。
% CHUNK appC1 DONE
